\chapter{注意力汇的计算动力学：归一化、结构角色与稀疏路由}
\label{chp:topology_of_attention_sink}

注意力汇提出了一个容易被几何隐喻遮蔽的问题：为什么因果语言模型会把相当一部分注意力分配给起始符、换行或其他内容贡献很弱的位置？这种现象既不能仅凭Softmax的和为一约束完整解释，也不能据此断言模型内部出现了虫洞、度量奇点或小世界拓扑。我们把研究对象限定为注意力层中可观测的权重、值向量贡献、残差更新与任务风险，区分概率质量的去向、信息实际写入的位置和跨层形成的功能角色。本章进一步讨论注意力汇与巨激活的关系，说明稀疏世界关系并不等于稠密参数矩阵中的每个非零数都是虚假语义边，并把小世界结构从“物理必然”改写为受任务、预算和路由规则约束的候选设计。由此，HSF-HD描述结构与活动如何联合演化；AgentOS只作为一种可审计的工程实例，而不是所有模型的唯一实现。

\section{从归一化约束到注意力汇：定义、机制与必要条件}
\label{sec:softmax_surface_tension}

我们先定义可检验对象。设某一因果自注意力头在时刻$i$可访问位置集合为$\mathcal{P}_i$，查询、键和值分别为$q_i,k_j,v_j\in\mathbb{R}^{d_h}$。在温度$\tau>0$且掩码$m_{ij}\in\{0,-\infty\}$给定时，注意力权重与输出为
\begin{equation}
 a_{ij}=\frac{\exp(q_i^{\mathsf T}k_j/\tau+m_{ij})}
 {\sum_{r\in\mathcal{P}_i}\exp(q_i^{\mathsf T}k_r/\tau+m_{ir})},
 \qquad
 o_i=\sum_{j\in\mathcal{P}_i}a_{ij}v_jW_O.
\end{equation}
其中$a_{ij}$是无量纲概率权重，$o_i$是与残差通道维度一致的向量，$W_O$为输出投影。给定位置$s$，我们把它称为该头在数据集$D$上的候选注意力汇，当且仅当跨越足够多的查询位置、序列与样本后，$a_{is}$持续高于位置匹配基线，而且这种集中不能完全由少数特定词义或短距离邻接解释。这个定义只是统计筛选；是否承担稳定计算角色，还要测量值贡献$\|a_{is}v_sW_O\|$、消融后的损失变化和跨层可重复性。

Softmax确实制造了一个局部约束：对每个查询都有$\sum_j a_{ij}=1$。当各个内容位置的得分都不值得写入时，权重仍必须落到某处。如果训练能够构造一个“高键匹配、低值贡献”的位置$s$，模型便可令$q_i^{\mathsf T}k_s$较高，同时令$\|v_sW_O\|$较小，把不用的概率质量放在那里而少改动残差流。这是注意力汇的一种充分机制，不是唯一机制。多头之间可能分工，一些头使用起始符作为空操作通道，另一些头把它当作文档身份或全局摘要；残差连接也允许注意力层输出接近零而主通道继续传播。因而“权重大”不能直接翻译为“携带信息多”，更不能翻译为活动强度、信息熵或设备能耗大。

仅有归一化仍不足以指定汇的位置。位置$s$必须在训练分布中稳定可见，其键表示必须容易被许多查询访问，优化过程还要发现把无用权重引向$s$比重塑所有内容得分更便宜。起始符满足固定身份和可用范围两个条件；在因果掩码下，它对后续所有位置均可见，而靠后的标记没有这种覆盖。标点、换行或填充模式也可能因高频、位置稳定或数据模板而获得类似角色。若模型改用允许“拒绝关注”的门控、显式零值槽、非归一化核注意力或局部窗口，汇的强度和位置可能改变。这些干预能区分Softmax压力、因果可达性、数据频率与优化偏好，而无须把Softmax称为表面张力。

我们还要区别训练形成与推理维持。设参数在离散训练步$k$按
\begin{equation}
 \theta_{k+1}=\theta_k-\eta_k\widehat\nabla_\theta
 \bigl[L_{D_k}(\theta_k)+\lambda_k\Omega(\theta_k)\bigr]
\end{equation}
更新，其中$\eta_k$为步长，$D_k$为批次，$L$和$\Omega$均为经过协议定义的无量纲目标。汇是数据、初始化、掩码、位置编码、优化器与正则共同塑造的训练结果；在推理时，给定输入只是在已有参数上产生一次条件化读出。若数据模板或上下文长度改变，汇可能迁移或分裂。损失下降说明训练目标改善，不证明该位置是全局参照系，也不证明长上下文任务已经成功。

验证协议至少包括四个对照。第一，保留序列内容而置换候选汇的位置，观察集中是否跟随身份还是绝对位置；第二，分别消融键、值与输出投影，区分概率吸收和有效写入；第三，比较Softmax、带显式空槽的门控和局部或线性注意力，检验归一化是否必要；第四，在长度外推、跨模板和跨领域样本上重复测量。只有权重统计、因果干预与任务指标共同支持，才能把某个位置称为功能性注意力汇。曲线相似、位置固定或单次可视化都不足以推出拓扑不变量。

还需把汇强度与最终输出分开报告。同一权重变化经不同值向量和后续层放大后可能产生完全不同的预测效应，因此我们同时保存查询位置、头身份、残差增量与输出差异。若只在平均权重上观察到集中，却没有稳定贡献或因果效应，我们只称它为分配模式，不把它升级为全局计算枢纽。

\section{结构枢纽、语义空载与巨激活：三种现象不能合并}
\label{sec:geometric_hubs}

把注意力汇直接称作小世界枢纽会混合三个层级。第一层是注意力矩阵中的高入权重位置，它描述特定层、特定头和特定输入下的活动分配；第二层是跨样本保持的计算角色，例如空操作、分隔、位置锚定或摘要；第三层才是把多次交互聚合后得到的功能网络中心性。设$A^{(\ell,h)}(x)$为样本$x$在第$\ell$层第$h$头的注意力矩阵，我们可以在声明聚合规则后定义
\begin{equation}
 C_j=\mathbb{E}_{x,\ell,h,i}\left[A^{(\ell,h)}_{ij}(x)\right],
\end{equation}
但$C_j$只是平均入权重，不等于图论中基于最短路径的介数中心性。注意力边随输入变化、带方向且跨层复合，不能从$C_j$较高推出“任意信息都经过该点”，更不能声称路径长度必然从$O(N)$降为$O(\log N)$。

\label{sec:semantic_vacuum_theorem}
所谓“语义真空”更适合改写为任务相关的低内容干扰。把候选汇的键和值分别记为$k_s$与$v_s$。高键可达性要求许多查询与$k_s$具有较高相似度，低写入要求投影后$\|v_sW_O\|$较小，二者由不同参数控制。若$v_s$近零，该位置可以吸收权重而不显著改变输出；若$v_s$编码文档类型、边界或全局状态，它仍可能是有意义的控制标记。我们因此不提出“枢纽必须语义为空”的定理，而给出可测假说：当某一头主要承担空操作时，提高汇位置值向量的内容相关变化会增加输出干扰；当该头承担摘要或广播时，强行压低值贡献反而会损害任务。两类结果都允许存在，角色必须由消融和替换实验识别。

“股票经由猫节点便被猫语义污染”的叠加故事也不构成推导。Transformer的输出是经学习投影、归一化、非线性和残差组合后的向量，某个值方向是否影响最终预测取决于后续读出灵敏度。我们用局部干预增益衡量干扰：对位置$s$施加有界扰动$\delta v_s$，定义
\begin{equation}
 G_s(\delta)=\mathbb{E}_{x\sim D}
 \frac{\|f_{\theta}(x;v_s+\delta v_s)-f_{\theta}(x;v_s)\|}
 {\|\delta v_s\|+\epsilon},
\end{equation}
其中输入、层和扰动空间都应固定，$\epsilon>0$避免零除。低$G_s$支持该位置在所测任务上近似空操作，高$G_s$则提示其携带控制或内容。这个量是有限扰动下的经验灵敏度，不是波函数污染或逻辑退相干。

\label{sec:metric_singularity_energy}
巨激活是另一个需要独立定义的现象。设隐藏状态$z_i^{(\ell)}\in\mathbb{R}^{d}$，我们可用相对于同层分布的稳健标准分数识别异常坐标，例如
\begin{equation}
 M_{i,r}^{(\ell)}=
 \frac{|z_{i,r}^{(\ell)}-\operatorname{median}(z_{:,r}^{(\ell)})|}
 {\operatorname{MAD}(z_{:,r}^{(\ell)})+\epsilon}.
\end{equation}
这里$M$是无量纲统计量；它不等于注意力权重，更不是应力—能量张量。大激活可能服务于归一化偏置、门控、条件控制或量化尺度，也可能只是训练和数值表示的副产物。由于Softmax权重始终不超过一，原稿中$A_{sink}\to\infty$与$d^2\propto-\log A$的组合在定义域上就不成立。即便某个打分logit很大，也不能由此推出黎曼距离趋零或时空曲率发散。

注意力汇与巨激活可以相关，却不存在定义上的同一性。我们需要在层、头、位置和训练检查点上建立联合时间序列，比较汇强度$C_s$、值贡献、异常激活$M$与任务风险。随后分别执行裁剪激活、重参数化归一化、替换候选汇和保持函数近似不变的尺度变换。如果激活被压缩而汇仍存在，巨激活不是必要条件；如果汇位置替换后激活迁移，二者可能共享控制角色；如果两者消失而任务不变，它们可能只是特定参数化的实现。这样的条件结论比“高曲率需要高能量”的类比更接近可重复机制。

反例同样重要：某些头可能没有稳定汇，某些汇携带强控制语义，某些巨激活也可能出现在普通内容位置。验证报告应保留这些负例，并说明阈值、样本和不确定性。只选择最醒目的头作图会造成幸存者偏差，也无法支持关于整个模型拓扑的普遍结论。

% 第三节续写位置：世界关系图与模型计算图；随后完成稀疏路由工程闭环。
